% !TEX encoding = UTF-8
% !TEX program = xelatex
\documentclass[11pt,a4paper]{article}

\usepackage[a4paper,margin=1.6cm]{geometry}
\usepackage{fontspec}
\usepackage{hyperref}
\usepackage{xcolor}
\usepackage{setspace}
\usepackage{tabularx}
\usepackage{graphicx}

\hypersetup{
  colorlinks=true,
  linkcolor=black,
  urlcolor=blue,
  citecolor=black
}

\defaultfontfeatures{Ligatures=TeX}
\setmainfont[
  BoldFont={Songti SC Bold},
  ItalicFont={Kaiti SC}
]{Songti SC}
\setsansfont{PingFang SC}
\newfontfamily\latinfont{Times New Roman}

\XeTeXlinebreaklocale "zh"
\XeTeXlinebreakskip = 0pt plus 1pt

\pagestyle{empty}
\setstretch{1.08}
\hyphenpenalty=10000
\exhyphenpenalty=10000
\emergencystretch=2em

\newcommand{\cvsection}[1]{%
  \par\vspace{0.75em}
  \noindent{\large\bfseries #1}\par\vspace{0.15em}\hrule\vspace{0.35em}
}

\newcommand{\cvitem}[2]{%
  \par\vspace{0.5em}
  \noindent
  \begin{tabularx}{\linewidth}{@{}X r@{}}%
    \textbf{#2} & \textbf{#1}\\%
  \end{tabularx}%
  \par\vspace{0.1em}
}

\newcommand{\cvsubitem}[1]{\par\vspace{0.02em}\noindent{\itshape #1}\\[-0.95em]}
\newenvironment{cvitemize}{%
  \begin{itemize}%
  \setlength{\leftmargin}{1.6em}%
  \setlength{\itemsep}{0pt}%
  \setlength{\topsep}{-0.15em}%
  \setlength{\parsep}{0pt}%
  \setlength{\parskip}{0pt}%
}{\end{itemize}\vspace{-0.25em}}

\newcommand{\cvphoto}{jp-1-revised.png}

\begin{document}
\vspace*{-1.0cm}

\noindent
\begin{tabularx}{\linewidth}{@{}X m{2.8cm}@{}}
  \begin{minipage}[t]{\linewidth}
    \vspace*{-5.8em}
    {\LARGE\bfseries 蒋鹏}\\[0.8em]
    \textbf{电话：}13161911011\\
    \textbf{邮箱：}\href{mailto:jiang-p21@mails.tsinghua.edu.cn}{jiang-p21@mails.tsinghua.edu.cn}\\
    \textbf{主页：}\href{https://jp-17.github.io/}{https://jp-17.github.io/}\\
  \end{minipage}
  &
  \centering\raisebox{3.8em}{\includegraphics[width=2.8cm]{\cvphoto}}
\end{tabularx}
\vspace*{-6.0em}

\cvsection{教育背景}
\begingroup
\renewcommand{\cvitem}[2]{%
  \par\vspace{0.35em}
  \noindent
  \begin{tabularx}{\linewidth}{@{}X r@{}}%
    #2 & \textbf{#1}\\%
  \end{tabularx}%
  \par\vspace{0.15em}
}
\cvitem{2017.08--2021.06}{清华大学\quad 生命科学学院\quad 学士学位}
\cvitem{2018.08--2021.06}{清华大学\quad 计算机科学与技术系\quad 本科辅修}
\cvitem{2021.08--现在}{清华大学\quad 生命科学学院\quad 计算神经科学博士生\quad 导师：贾晓轩}
\endgroup

\cvsection{研究方向}
\noindent\textbf{关键词：}Neural Foundation Model，Neural Encoding \& Decoding，World Model，Neural Representation \\[0.25em]
\noindent\textbf{当前研究：}聚焦 neural foundation model，探索如何构建神经预训练模型进行神经编码与解码的方向研究。

\cvsection{学术经历}
\cvitem{2022--2024}{多任务学习的神经表征分析研究}
\cvsubitem{清华大学\quad 生命科学学院\quad 贾晓轩实验室}
\begin{cvitemize}
  \item 借助 PCA/CCA 等表征降维方式，对不同视觉脑区神经活动在低维空间进行可视化，分析小鼠视觉编码的神经动力学特征与脑区间信息传递方式；
  \item 构建 RNN 模型模拟不同脑区完成多种认知任务，并通过 fixed point analysis 来研究网络本身的动力学表征及其完成任务的机制；
  \item 微调训练大语言模型进行多种认知任务的持续学习，通过分析大模型本身的权重和神经元激活特征的稀疏分解，来研究任务的关系是否会呈现出组合泛化的表征形式。
\end{cvitemize}

\cvitem{2026}{NeuroHorizon: Long-Horizon Forward Prediction of Neural Population Activity via Autoregressive Decoding with Hierarchical Memory}
\cvsubitem{Peng Jiang, Xiaoxuan Jia\textsuperscript{*} · Submitted to NeurIPS 2026 · Under review}
\begin{cvitemize}
  \item 开发神经活动预测模型，尽量减少模型在神经活动的长时程前向预测上的 rollout prediction 性能衰减;
  \item 模型利用 POYO-like spike encoding model 进行高分辨率的神经信息编码， 同时在 decoder 侧结合 tail+segment 的分层记忆机制的减少预测性能在长时程情况下的衰减;
  \item 在多个数据集上评估模型的长时程多步 rollout 预测神经活动的稳定性，以及模型的 scaling \& generalization 的能力。
\end{cvitemize}

\cvsection{项目经历}
\cvitem{2024--2025}{基于 3DGS 的实时音频驱动数字人}
\cvsubitem{空间漫步（深圳）科技有限公司\quad 联合创始人-算法负责(休学创业)}
\begin{cvitemize}
  \item 基于单目视频进行人物 3D 模型重建的优化，包括 flame tracking 和 3dgs model 的优化；
  \item 改进基于 diffusion model 为框架的音频驱动模型，提高驱动模型的生动性和准确性。
\end{cvitemize}

\cvitem{2023.06}{NeuroGPT 项目}
\cvsubitem{清华大学大模型应用设计挑战赛 · 第4名/共94支队伍}
\begin{cvitemize}
  \item 以多模态大语言模型作为信息的中转，通过脑电解码语义后再映射解码为动作/操控等信息，从而赋能医疗/教育/娱乐等广泛场景。
\end{cvitemize}

\end{document}
